\section{Introduction}

As AI systems scale in capability and autonomy, their most consequential failures increasingly arise not from insufficient intelligence, but from \emph{unexamined reasoning paths}. Plans that are internally coherent yet externally flawed can pass unnoticed through single-agent pipelines, especially when systems are optimized for fluency, speed, or apparent confidence rather than adversarial scrutiny.

Current mitigation strategies typically operate along one of three axes: (1) alignment objectives and reward modeling, (2) self-consistency or internal verification techniques, and (3) post-hoc evaluation and monitoring. While valuable, these approaches share a common limitation: they do not structurally require disagreement prior to execution.

In complex engineered systems, reliability is rarely achieved through correctness alone. Instead, it emerges from organizational separation of concerns, explicit review authority, and bounded escalation paths. Safety-critical domains rely on independent checks, veto power, and audit trails—not trust in a single component’s internal reasoning.

This paper applies that insight to AI orchestration.

We introduce the \emph{Rival Orchestration Kernel (ROK)}, a minimal, open reference architecture that enforces adversarial review before execution through explicit role separation and protocolized interaction. ROK is not a prompting strategy, a training method, or an alignment framework. It is an execution-time reliability mechanism that treats disagreement as a first-class primitive.

ROK decomposes a task into four roles: a Planner, a Critic, a Decision layer, and an Executor. These roles interact through a bounded revise–recheck loop. If a plan is vetoed, revision is permitted only a finite number of times, after which execution is either halted or forced via an explicit override. All steps emit structured, schema-versioned trace events in an append-only log, enabling replay, validation, and downstream analysis.

ROK does not claim to ensure correctness of outputs, nor does it replace domain expertise or human judgment. Instead, it provides a substrate for failure interception and visibility, shifting error detection earlier in the execution pipeline—before irreversible actions occur.

